\chapter{The Basics}

This chapter develops the basic order-theoretic algebra that
underpins the rest of the library: lattice-ordered abelian groups
and the identities governing positive part, negative part and
absolute value; vector lattices and their interaction with scalar
multiplication; the Archimedean property; the Riesz decomposition
theorem; normed and Banach lattices and the automatic continuity of
the lattice operations; the calculus of disjointness; and finally
the two order-completeness conditions (Dedekind and sigma Dedekind)
together with their restriction to the positive cone.

\section{Vector lattices and the Archimedean property}

This section sets up the basic order-theoretic algebra underpinning the
rest of the library. We work first in the generality of a lattice-ordered
abelian group, establishing the fundamental identities for the positive
part $x^+$, the negative part $x^-$, and the absolute value $\lvert x\rvert$,
together with the uniqueness of the positive part and the interaction of
the lattice operations with translation and with arbitrary suprema and
infima. We then specialise to \emph{vector lattices}, i.e.\ real vector
spaces with a compatible lattice structure, and record how non-negative
scalars distribute over the lattice operations, culminating in the identity
$\lvert \lambda x\rvert = \lvert \lambda\rvert\,\lvert x\rvert$. The section
closes with the Archimedean property, which plays a foundational role in
later chapters: we give two equivalent formulations and use them to show
that an Archimedean vector lattice of dimension strictly greater than
$1$ necessarily contains vectors that are neither non-negative nor
non-positive.

\subsection{Lattice-ordered abelian groups}

Throughout this subsection $X$ denotes a \emph{lattice-ordered abelian
group}: an additive abelian group equipped with a lattice order whose
translations are order-preserving. We write
$x^+ := x \vee 0$, $x^- := (-x) \vee 0$ and
$\lvert x \rvert := x \vee (-x)$ for the positive part, negative part
and absolute value of $x \in X$.

\begin{lemma}
\label{lem:abs-eq-zero-iff-zero}
\lean{abs_eq_zero_iff_zero}
\leanok
Let $X$ be a lattice-ordered abelian group and let $x \in X$. Then
$\lvert x \rvert = 0$ if and only if $x = 0$.
\end{lemma}

\begin{lemma}
\label{lem:uniqueness-posPart}
\lean{uniqueness_posPart}
\leanok
Let $X$ be a lattice-ordered abelian group, let $x \in X$, and suppose
$u, v \in X$ satisfy $x = u - v$ and $u \wedge v = 0$. Then $u = x^+$.
\end{lemma}

\begin{lemma}
\label{lem:sup-eq-add-posPart}
\lean{sup_eq_add_posPart}
\leanok
For all $x, y$ in a lattice-ordered abelian group,
$x \vee y = x + (y - x)^+.$
\end{lemma}

\begin{lemma}
\label{lem:inf-eq-sub-posPart}
\lean{inf_eq_sub_posPart}
\leanok
For all $x, y$ in a lattice-ordered abelian group,
$x \wedge y = x - (x - y)^+.$
\end{lemma}

\begin{lemma}
\label{lem:sub-inf-eq-posPart}
\lean{sub_inf_eq_posPart}
\leanok
For all $x, y$ in a lattice-ordered abelian group,
$x - x \wedge y = (x - y)^+.$
\end{lemma}

\begin{lemma}
\label{lem:posPart-add-le}
\lean{posPart_add_le}
\leanok
For all $x, y$ in a lattice-ordered abelian group,
$(x + y)^+ \le x^+ + y^+.$
\end{lemma}

\begin{lemma}
\label{lem:inf-le-inf-add-inf-of-nonneg}
\lean{inf_le_inf_add_inf_of_nonneg}
\leanok
Let $X$ be a lattice-ordered abelian group and let $x, a, b \in X$ with
$0 \le x$, $0 \le a$, and $0 \le b$. Then
$x \wedge (a + b) \le x \wedge a + x \wedge b.$
\end{lemma}

\begin{lemma}
\label{lem:isLUB-const-add}
\lean{isLUB_const_add}
\leanok
Let $X$ be a lattice-ordered abelian group, let $A \subseteq X$ and
$a, x \in X$. If $a = \sup A$, then
$x + a = \sup \{x + z : z \in A\}$.
\end{lemma}

\begin{lemma}
\label{lem:isGLB-const-add}
\lean{isGLB_const_add}
\leanok
Let $X$ be a lattice-ordered abelian group, let $A \subseteq X$ and
$a, x \in X$. If $a = \inf A$, then
$x + a = \inf \{x + z : z \in A\}$.
\end{lemma}

\begin{lemma}
\label{lem:isLUB-inf-const}
\lean{isLUB_inf_const}
\leanok
\uses{lem:inf-eq-sub-posPart, lem:sub-inf-eq-posPart}
Let $X$ be a lattice-ordered abelian group, let $A \subseteq X$ and
$a, x \in X$. If $a = \sup A$, then
$x \wedge a = \sup \{x \wedge z : z \in A\}$.
\end{lemma}

\begin{lemma}
\label{lem:isGLB-sup-const}
\lean{isGLB_sup_const}
\leanok
\uses{lem:sup-eq-add-posPart}
Let $X$ be a lattice-ordered abelian group, let $A \subseteq X$ and
$a, x \in X$. If $a = \inf A$, then
$x \vee a = \inf \{x \vee z : z \in A\}$.
\end{lemma}

\begin{lemma}
\label{lem:isGLB-add-sets}
\lean{isGLB_add_sets}
\leanok
Let $X$ be a lattice-ordered abelian group, let $A, B \subseteq X$ and
$a, b \in X$. If $a = \inf A$ and $b = \inf B$, then
$a + b = \inf \{p + q : p \in A,\ q \in B\}$.
\end{lemma}

\begin{lemma}
\label{lem:isGLB-sup-sets}
\lean{isGLB_sup_sets}
\leanok
\uses{lem:isGLB-sup-const}
Let $X$ be a lattice-ordered abelian group, let $A, B \subseteq X$ and
$a, b \in X$. If $a = \inf A$ and $b = \inf B$, then
$a \vee b = \inf \{p \vee q : p \in A,\ q \in B\}$.
\end{lemma}

\begin{lemma}
\label{lem:isGLB-union}
\lean{isGLB_union}
\leanok
Let $X$ be a lattice, let $A, B \subseteq X$ and $a, b \in X$. If
$a = \inf A$ and $b = \inf B$, then $a \wedge b = \inf (A \cup B)$.
\end{lemma}

\subsection{Vector lattices}

A \emph{vector lattice} is a lattice-ordered abelian group equipped with
a real vector space structure whose scalar multiplication by non-negative
reals is monotone. Throughout this subsection $X$ denotes a vector
lattice.

\begin{definition}
\label{def:vector-lattice}
\lean{VectorLattice}
\leanok
Let $X$ be an additive abelian group equipped with a lattice structure
compatible with addition. A \emph{vector lattice} structure on $X$
consists of a real vector space structure on $X$ for which scalar
multiplication by non-negative reals is monotone: for all $\lambda \ge 0$
and $x \le y$ in $X$, $\lambda x \le \lambda y$.
\end{definition}

\begin{lemma}
\label{lem:nonneg-smul-sup}
\lean{nonneg_smul_sup}
\leanok
Let $X$ be a vector lattice, let $x, y \in X$ and let $\lambda \in
\mathbb{R}$ with $\lambda \ge 0$. Then
$\lambda \cdot (x \vee y) = (\lambda x) \vee (\lambda y).$
\end{lemma}

\begin{lemma}
\label{lem:isLUB-smul-of-nonneg}
\lean{isLUB_smul_of_nonneg}
\leanok
Let $X$ be a vector lattice, let $A \subseteq X$ and $a \in X$, and let
$\lambda \in \mathbb{R}$ with $\lambda \ge 0$. If $a = \sup A$, then
$\lambda a = \sup \{\lambda z : z \in A\}$.
\end{lemma}

\begin{lemma}
\label{lem:isGLB-smul-of-nonneg}
\lean{isGLB_smul_of_nonneg}
\leanok
Let $X$ be a vector lattice, let $A \subseteq X$ and $a \in X$, and let
$\lambda \in \mathbb{R}$ with $\lambda \ge 0$. If $a = \inf A$, then
$\lambda a = \inf \{\lambda z : z \in A\}$.
\end{lemma}

\begin{lemma}
\label{lem:nonneg-smul-inf}
\lean{nonneg_smul_inf}
\leanok
\uses{lem:nonneg-smul-sup}
Let $X$ be a vector lattice, let $x, y \in X$ and let $\lambda \in
\mathbb{R}$ with $\lambda \ge 0$. Then
$\lambda \cdot (x \wedge y) = (\lambda x) \wedge (\lambda y).$
\end{lemma}

\begin{lemma}
\label{lem:sup-smul-nonneg}
\lean{sup_smul_nonneg}
\leanok
Let $X$ be a vector lattice, let $x \in X$ with $0 \le x$ and let
$\alpha, \beta \in \mathbb{R}$. Then
$(\alpha \vee \beta) \cdot x = (\alpha x) \vee (\beta x).$
\end{lemma}

\begin{lemma}
\label{lem:inf-smul-nonneg}
\lean{inf_smul_nonneg}
\leanok
Let $X$ be a vector lattice, let $x \in X$ with $0 \le x$ and let
$\alpha, \beta \in \mathbb{R}$. Then
$(\alpha \wedge \beta) \cdot x = (\alpha x) \wedge (\beta x).$
\end{lemma}

\begin{theorem}
\label{thm:abs-smul}
\lean{abs_smul'}
\leanok
\uses{lem:nonneg-smul-sup}
Let $X$ be a vector lattice, let $x \in X$ and $\alpha \in \mathbb{R}$.
Then $\lvert \alpha x \rvert = \lvert \alpha\rvert \cdot \lvert x\rvert.$
\end{theorem}

\begin{lemma}
\label{lem:disjoint-smul}
\lean{disjoint_smul}
\leanok
\uses{lem:nonneg-smul-inf}
Let $X$ be a vector lattice, let $x, y \in X$ with $x \wedge y = 0$ and
let $\lambda \in \mathbb{R}$ with $0 \le \lambda$. Then
$(\lambda x) \wedge y = 0.$
\end{lemma}

\subsection{The Archimedean property}

A lattice-ordered abelian group is \emph{Archimedean} when no non-zero
positive element is infinitesimal with respect to some other element.
We record two equivalent formulations of this property and deduce that
in an Archimedean vector lattice of dimension at least two there exist
vectors which are neither negative nor positive.

\begin{definition}
\label{def:is-vl-archimedean}
\lean{IsVLArchimedean}
\leanok
A lattice-ordered abelian group $X$ is called \emph{Archimedean} when,
for every $x, y \in X$ with $0 \le x$ and $n \cdot x \le y$ for all
$n \in \mathbb{N}$, one has $x = 0$.
\end{definition}

\begin{lemma}
\label{lem:is-vl-archimedean-iff-le-zero}
\lean{isVLArchimedean_iff_le_zero_of_forall_nsmul_le}
\leanok
\uses{def:is-vl-archimedean}
Let $X$ be a lattice-ordered abelian group. Then $X$ is Archimedean if
and only if for all $x, y \in X$ the condition $n \cdot x \le y$ for all
$n \in \mathbb{N}$ implies $x \le 0$.
\end{lemma}

\begin{lemma}
\label{lem:is-vl-archimedean-iff-isGLB-inv-smul}
\lean{isVLArchimedean_iff_isGLB_inv_smul}
\leanok
\uses{def:is-vl-archimedean, lem:is-vl-archimedean-iff-le-zero}
A vector lattice $X$ is Archimedean if and only if for every $u \in X$
with $0 \le u$,
$\inf \{ n^{-1} \cdot u : n \in \mathbb{N},\ n > 0 \} = 0$.
\end{lemma}

\begin{lemma}
\label{lem:infinitesimal-eq-zero}
\lean{infinitesimal_eq_zero}
\leanok
\uses{def:is-vl-archimedean, lem:abs-eq-zero-iff-zero}
Let $X$ be an Archimedean lattice-ordered abelian group and let
$x, y \in X$. If $n \cdot \lvert x \rvert \le y$ for all $n \in
\mathbb{N}$, then $x = 0$.
\end{lemma}

\begin{theorem}
\label{thm:exists-not-nonneg-not-nonpos}
\lean{exists_not_nonneg_not_nonpos_of_one_lt_rank}
\leanok
\uses{def:is-vl-archimedean, lem:infinitesimal-eq-zero, thm:abs-smul}
Let $X$ be an Archimedean vector lattice with dimension
strictly greater than $1$. Then there exists $x \in X$ which is neither
negative nor positive.
\end{theorem}

\section{Riesz decomposition}

This section collects the Riesz decomposition theorem and its standard
refinements in a lattice-ordered abelian group $X$. The binary form
splits a positive element dominated by a sum of two positive elements
into a corresponding sum. Iterating and combining the argument yields a
matrix form for two partitions of the same positive sum, a version with
a finite sum on the left hand side, a decomposition adapted to the
absolute value, and an identification of sum intervals
$[0, x] + [0, y] = [0, x + y]$ for $x, y \ge 0$.

\begin{theorem}
\label{thm:riesz-decomposition}
\lean{riesz_decomposition}
\leanok
Let $X$ be a lattice-ordered abelian group and let $x, y, z \in X$ with
$0 \le x$, $0 \le y$, $0 \le z$ and $x \le y + z$. Then there exist
$x_1, x_2 \in X$ with $0 \le x_1 \le y$, $0 \le x_2 \le z$ and
$x = x_1 + x_2$.
\end{theorem}

\begin{theorem}
\label{thm:riesz-decomposition-matrix}
\lean{riesz_decomposition_matrix}
\leanok
\uses{thm:riesz-decomposition}
Let $X$ be a lattice-ordered abelian group, let $m, n \in \mathbb{N}$,
let $u : \{0, \dots, m-1\} \to X$ and $z : \{0, \dots, n-1\} \to X$
with $0 \le u_i$ for all $i$ and $0 \le z_j$ for all $j$, and
$\sum_i u_i = \sum_j z_j$. Then there exists a family
$x : \{0, \dots, m-1\} \times \{0, \dots, n-1\} \to X$ with $0 \le x_{ij}$
for all $i, j$ such that $u_i = \sum_j x_{ij}$ for every $i$ and
$z_j = \sum_i x_{ij}$ for every $j$.
\end{theorem}

\begin{theorem}
\label{thm:riesz-decomposition-sum-le}
\lean{riesz_decomposition_sum_le}
\leanok
\uses{thm:riesz-decomposition, thm:riesz-decomposition-matrix}
Let $X$ be a lattice-ordered abelian group, let $k \in \mathbb{N}$, let
$u : \{0, \dots, k-1\} \to X$ and $x, y \in X$ with $0 \le u_i$ for all
$i$, $0 \le x$, $0 \le y$ and $\sum_i u_i \le x + y$. Then there exist
$v, w : \{0, \dots, k-1\} \to X$ with $0 \le v_i$ and $0 \le w_i$ for
all $i$, $\sum_i v_i \le x$, $\sum_i w_i \le y$, and $u_i = v_i + w_i$
for every $i$.
\end{theorem}

\begin{theorem}
\label{thm:exists-abs-decomposition}
\lean{exists_abs_decomposition}
\leanok
\uses{thm:riesz-decomposition, lem:uniqueness-posPart}
Let $X$ be a lattice-ordered abelian group, let $x, u, v \in X$ with
$0 \le u$, $0 \le v$ and $\lvert x \rvert = u + v$. Then there exist
$y, z \in X$ with $x = y + z$, $\lvert y \rvert = u$ and
$\lvert z \rvert = v$.
\end{theorem}

\begin{theorem}
\label{thm:riesz-decomposition-abs}
\lean{riesz_decomposition_abs}
\leanok
\uses{thm:riesz-decomposition, thm:exists-abs-decomposition}
Let $X$ be a lattice-ordered abelian group, let $n \in \mathbb{N}$, let
$x \in X$ and let $v : \{0, \dots, n-1\} \to X$ with $0 \le v_i$ for
every $i$ and $\lvert x \rvert \le \sum_i v_i$. Then there exists
$y : \{0, \dots, n-1\} \to X$ with $x = \sum_i y_i$ and
$\lvert y_i \rvert \le \lvert x \rvert \wedge v_i$ for every $i$.
\end{theorem}

\begin{theorem}
\label{thm:Icc-add-Icc-of-nonneg}
\lean{Icc_add_Icc_of_nonneg}
\leanok
\uses{thm:riesz-decomposition}
Let $X$ be a lattice-ordered abelian group and let $x, y \in X$ with
$0 \le x$ and $0 \le y$. Then
$[0, x] + [0, y] = [0, x + y]$
as subsets of $X$.
\end{theorem}

\section{Normed and Banach lattices}

A \emph{normed vector lattice} is a real vector lattice equipped with a
norm satisfying the solid axiom $\lvert x \rvert \le \lvert y \rvert
\Rightarrow \lVert x \rVert \le \lVert y \rVert$. This compatibility
between the norm and the lattice operations has a number of immediate
consequences: the lattice operations $\vee$, $\wedge$ and
$\lvert \cdot \rvert$ are norm-continuous, the space is automatically
Archimedean, the positive cone is norm-closed, order intervals are
norm-closed and norm-bounded, and monotone norm-convergent sequences
converge to their supremum. A \emph{Banach lattice} is a normed vector
lattice whose norm is complete.

\begin{definition}
\label{def:normed-vector-lattice}
\lean{NormedVectorLattice}
\leanok
\uses{def:vector-lattice}
Let $X$ be a normed abelian group equipped with a compatible lattice
structure. A \emph{normed vector lattice} structure on $X$ consists of
a vector lattice structure together with the solid axiom
$\lvert x \rvert \le \lvert y \rvert \implies \lVert x \rVert \le
\lVert y \rVert$ for all $x, y \in X$, and the identity $\lVert
\lambda x \rVert = \lvert \lambda \rvert \, \lVert x \rVert$ for all
$\lambda \in \mathbb{R}$ and $x \in X$.
\end{definition}

\begin{lemma}
\label{lem:normed-vector-lattice-normed-space}
\lean{NormedVectorLattice.instNormedSpace}
\leanok
\uses{def:normed-vector-lattice}
Every normed vector lattice is a normed space over $\mathbb{R}$.
\end{lemma}

\begin{theorem}
\label{thm:continuous-sup}
\lean{NormedVectorLattice.continuous_sup}
\leanok
\uses{def:normed-vector-lattice}
Let $X$ be a normed vector lattice. The map $(x, y) \mapsto x \vee y$
from $X \times X$ to $X$ is continuous.
\end{theorem}

\begin{theorem}
\label{thm:continuous-inf}
\lean{NormedVectorLattice.continuous_inf}
\leanok
\uses{def:normed-vector-lattice}
Let $X$ be a normed vector lattice. The map $(x, y) \mapsto x \wedge y$
from $X \times X$ to $X$ is continuous.
\end{theorem}

\begin{theorem}
\label{thm:lipschitz-abs}
\lean{NormedVectorLattice.lipschitzWith_abs}
\leanok
\uses{def:normed-vector-lattice}
Let $X$ be a normed vector lattice. The map $x \mapsto \lvert x \rvert$
from $X$ to $X$ is Lipschitz with constant $1$.
\end{theorem}

\begin{theorem}
\label{thm:normed-vector-lattice-archimedean}
\lean{NormedVectorLattice.instIsVLArchimedean}
\leanok
\uses{def:normed-vector-lattice, def:is-vl-archimedean}
Every normed vector lattice is Archimedean.
\end{theorem}

\begin{lemma}
\label{lem:normed-vector-lattice-orderClosedTopology}
\lean{NormedVectorLattice.instOrderClosedTopology}
\leanok
\uses{def:normed-vector-lattice}
Let $X$ be a normed vector lattice. The order relation on $X$ is closed
in the norm topology.
\end{lemma}

\begin{theorem}
\label{thm:isClosed-nonneg-cone}
\lean{NormedVectorLattice.isClosed_nonneg_cone}
\leanok
\uses{lem:normed-vector-lattice-orderClosedTopology}
Let $X$ be a normed vector lattice. The positive cone
$\{x \in X : 0 \le x\}$ is norm-closed.
\end{theorem}

\begin{theorem}
\label{thm:le-of-tendsto-of-tendsto}
\lean{NormedVectorLattice.le_of_tendsto_of_tendsto}
\leanok
\uses{lem:normed-vector-lattice-orderClosedTopology}
Let $X$ be a normed vector lattice, let $u, v : \mathbb{N} \to X$ and
$a, b \in X$. If $u_n \to a$ and $v_n \to b$ in norm and
$u_n \le v_n$ for every $n$, then $a \le b$.
\end{theorem}

\begin{theorem}
\label{thm:isLUB-of-monotone-tendsto}
\lean{NormedVectorLattice.isLUB_of_monotone_tendsto}
\leanok
\uses{lem:normed-vector-lattice-orderClosedTopology}
Let $X$ be a normed vector lattice, let $u : \mathbb{N} \to X$ be
monotone and let $l \in X$. If $u_n \to l$ in norm, then $l$ is the
least upper bound of the range of $u$.
\end{theorem}

\begin{theorem}
\label{thm:isGLB-of-antitone-tendsto}
\lean{NormedVectorLattice.isGLB_of_antitone_tendsto}
\leanok
\uses{lem:normed-vector-lattice-orderClosedTopology}
Let $X$ be a normed vector lattice, let $u : \mathbb{N} \to X$ be
antitone and let $l \in X$. If $u_n \to l$ in norm, then $l$ is the
greatest lower bound of the range of $u$.
\end{theorem}

\begin{lemma}
\label{lem:isClosed-interval}
\lean{NormedVectorLattice.isClosed_interval}
\leanok
\uses{lem:normed-vector-lattice-orderClosedTopology}
Let $X$ be a normed vector lattice and let $a, b \in X$. The order
interval $[a, b]$ is norm-closed.
\end{lemma}

\begin{lemma}
\label{lem:norm-le-norm-abs-sup-abs-of-mem-Icc}
\lean{NormedVectorLattice.norm_le_norm_abs_sup_abs_of_mem_Icc}
\leanok
\uses{def:normed-vector-lattice}
Let $X$ be a normed vector lattice, let $a, b \in X$ and let
$x \in [a, b]$. Then $\lVert x \rVert \le \lVert \lvert a \rvert \vee
\lvert b \rvert \rVert$.
\end{lemma}

\begin{theorem}
\label{thm:isBounded-of-bddBelow-bddAbove}
\lean{NormedVectorLattice.isBounded_of_bddBelow_bddAbove}
\leanok
\uses{lem:norm-le-norm-abs-sup-abs-of-mem-Icc}
Let $X$ be a normed vector lattice and let $s \subseteq X$ be both
bounded below and bounded above in the order. Then $s$ is norm-bounded.
\end{theorem}

\begin{definition}
\label{def:banach-lattice}
\lean{BanachLattice}
\leanok
\uses{def:normed-vector-lattice}
A \emph{Banach lattice} is a normed vector lattice whose underlying
normed space is complete.
\end{definition}

\section{Disjointness}

Two elements of a lattice-ordered abelian group are \emph{disjoint} when
the infimum of their absolute values vanishes. This section develops the
elementary calculus of disjointness: symmetry, the zero rules, the
Birkhoff identity $\lvert x + y \rvert = \lvert x \rvert + \lvert y \rvert$
for disjoint pairs, uniqueness of the positive/negative decomposition,
compatibility with scalar multiplication and with monotonicity under the
absolute value, and closure under finite suprema and sums. From these one
deduces the finite-family identity
$\lvert \sum_i \alpha_i x_i \rvert = \sum_i \lvert \alpha_i \rvert
\lvert x_i \rvert$ for pairwise-disjoint families, from which any
pairwise-disjoint family of non-zero vectors is linearly independent over
$\mathbb{R}$. The section closes with a normed-lattice statement: a norm
limit of a pairwise-disjoint sequence must be zero.

\begin{definition}
\label{def:isVLDisjoint}
\lean{IsVLDisjoint}
\leanok
Let $X$ be a lattice-ordered abelian group and let $x, y \in X$. We say
$x$ and $y$ are \emph{disjoint} when $\lvert x \rvert \wedge \lvert y
\rvert = 0$.
\end{definition}

\begin{lemma}
\label{lem:isVLDisjoint-comm}
\lean{isVLDisjoint_comm}
\leanok
\uses{def:isVLDisjoint}
Let $X$ be a lattice-ordered abelian group and let $x, y \in X$. Then
$x$ and $y$ are disjoint if and only if $y$ and $x$ are disjoint.
\end{lemma}

\begin{lemma}
\label{lem:isVLDisjoint-zero-left}
\lean{isVLDisjoint_zero_left}
\leanok
\uses{def:isVLDisjoint}
Let $X$ be a lattice-ordered abelian group and let $x \in X$. Then $0$
and $x$ are disjoint.
\end{lemma}

\begin{lemma}
\label{lem:isVLDisjoint-zero-right}
\lean{isVLDisjoint_zero_right}
\leanok
\uses{def:isVLDisjoint}
Let $X$ be a lattice-ordered abelian group and let $x \in X$. Then $x$
and $0$ are disjoint.
\end{lemma}

\begin{lemma}
\label{lem:inf-eq-zero-of-isVLDisjoint}
\lean{inf_eq_zero_of_isVLDisjoint}
\leanok
\uses{def:isVLDisjoint}
Let $X$ be a lattice-ordered abelian group and let $x, y \in X$ with
$0 \le x$ and $0 \le y$. If $x$ and $y$ are disjoint, then $x \wedge y
= 0$.
\end{lemma}

\begin{lemma}
\label{lem:isVLDisjoint-of-inf-eq-zero}
\lean{isVLDisjoint_of_inf_eq_zero}
\leanok
\uses{def:isVLDisjoint}
Let $X$ be a lattice-ordered abelian group and let $x, y \in X$ with
$x \wedge y = 0$. Then $x$ and $y$ are disjoint.
\end{lemma}

\begin{theorem}
\label{thm:isVLDisjoint-decomposition-unique}
\lean{isVLDisjoint_decomposition_unique}
\leanok
\uses{def:isVLDisjoint, lem:inf-eq-zero-of-isVLDisjoint, lem:uniqueness-posPart}
Let $X$ be a lattice-ordered abelian group and let $u_1, v_1, u_2, v_2
\in X$ be non-negative elements. If $u_1$ and $v_1$ are disjoint, $u_2$
and $v_2$ are disjoint, and $u_1 - v_1 = u_2 - v_2$, then $u_1 = u_2$
and $v_1 = v_2$.
\end{theorem}

\begin{theorem}
\label{thm:abs-add-of-isVLDisjoint}
\lean{abs_add_of_isVLDisjoint}
\leanok
\uses{def:isVLDisjoint, lem:uniqueness-posPart, lem:inf-le-inf-add-inf-of-nonneg}
Let $X$ be a lattice-ordered abelian group and let $x, y \in X$ be
disjoint. Then $\lvert x + y \rvert = \lvert x \rvert + \lvert y \rvert$.
\end{theorem}

\begin{lemma}
\label{lem:abs-sup-of-isVLDisjoint}
\lean{abs_sup_of_isVLDisjoint}
\leanok
\uses{def:isVLDisjoint}
Let $X$ be a lattice-ordered abelian group and let $x, y \in X$ with
$0 \le x$, $0 \le y$ and $x$, $y$ disjoint. Then $\lvert x \vee y
\rvert = \lvert x \rvert \vee \lvert y \rvert$.
\end{lemma}

\begin{lemma}
\label{lem:isVLDisjoint-posPart-negPart}
\lean{isVLDisjoint_posPart_negPart}
\leanok
\uses{def:isVLDisjoint, lem:isVLDisjoint-of-inf-eq-zero}
Let $X$ be a lattice-ordered abelian group and let $x \in X$. Then
$x^+$ and $x^-$ are disjoint.
\end{lemma}

\begin{lemma}
\label{lem:isVLDisjoint-mono-right}
\lean{IsVLDisjoint.mono_right}
\leanok
\uses{def:isVLDisjoint}
Let $X$ be a lattice-ordered abelian group and let $x, y, z \in X$. If
$x$ and $y$ are disjoint and $\lvert z \rvert \le \lvert y \rvert$, then
$x$ and $z$ are disjoint.
\end{lemma}

\begin{lemma}
\label{lem:isVLDisjoint-mono-left}
\lean{IsVLDisjoint.mono_left}
\leanok
\uses{def:isVLDisjoint, lem:isVLDisjoint-mono-right, lem:isVLDisjoint-comm}
Let $X$ be a lattice-ordered abelian group and let $x, y, z \in X$. If
$x$ and $y$ are disjoint and $\lvert z \rvert \le \lvert x \rvert$, then
$z$ and $y$ are disjoint.
\end{lemma}

\begin{lemma}
\label{lem:isVLDisjoint-smul-left}
\lean{IsVLDisjoint.smul_left}
\leanok
\uses{def:isVLDisjoint, thm:abs-smul, lem:disjoint-smul}
Let $X$ be a vector lattice and let $x, y \in X$ be disjoint. For every
$\alpha \in \mathbb{R}$, $\alpha x$ and $y$ are disjoint.
\end{lemma}

\begin{lemma}
\label{lem:isVLDisjoint-smul-right}
\lean{IsVLDisjoint.smul_right}
\leanok
\uses{def:isVLDisjoint, lem:isVLDisjoint-smul-left, lem:isVLDisjoint-comm}
Let $X$ be a vector lattice and let $x, y \in X$ be disjoint. For every
$\alpha \in \mathbb{R}$, $x$ and $\alpha y$ are disjoint.
\end{lemma}

\begin{theorem}
\label{thm:exists-pair-ne-zero-isVLDisjoint}
\lean{exists_pair_ne_zero_isVLDisjoint}
\leanok
\uses{def:isVLDisjoint, lem:isVLDisjoint-posPart-negPart, thm:exists-not-nonneg-not-nonpos}
Let $X$ be an Archimedean vector lattice with dimension strictly greater
than $1$. Then there exist non-zero $u, v \in X$ with $u$ and $v$
disjoint.
\end{theorem}

\begin{lemma}
\label{lem:add-eq-sup-of-isVLDisjoint-of-nonneg}
\lean{add_eq_sup_of_isVLDisjoint_of_nonneg}
\leanok
\uses{def:isVLDisjoint, lem:inf-eq-zero-of-isVLDisjoint}
Let $X$ be a lattice-ordered abelian group and let $x, y \in X$ with
$0 \le x$, $0 \le y$ and $x$, $y$ disjoint. Then $x + y = x \vee y$.
\end{lemma}

\begin{lemma}
\label{lem:sup-abs-eq-add-abs-of-isVLDisjoint}
\lean{sup_abs_eq_add_abs_of_isVLDisjoint}
\leanok
\uses{def:isVLDisjoint}
Let $X$ be a lattice-ordered abelian group and let $x, y \in X$ be
disjoint. Then $\lvert x \rvert \vee \lvert y \rvert = \lvert x \rvert
+ \lvert y \rvert$.
\end{lemma}

\begin{lemma}
\label{lem:abs-add-eq-sup-abs-of-isVLDisjoint}
\lean{abs_add_eq_sup_abs_of_isVLDisjoint}
\leanok
\uses{thm:abs-add-of-isVLDisjoint, lem:sup-abs-eq-add-abs-of-isVLDisjoint}
Let $X$ be a lattice-ordered abelian group and let $x, y \in X$ be
disjoint. Then $\lvert x + y \rvert = \lvert x \rvert \vee \lvert y
\rvert$.
\end{lemma}

\begin{lemma}
\label{lem:isVLDisjoint-sup-right}
\lean{IsVLDisjoint.sup_right}
\leanok
\uses{def:isVLDisjoint, lem:inf-le-inf-add-inf-of-nonneg}
Let $X$ be a lattice-ordered abelian group and let $x, y, z \in X$. If
$x$ and $y$ are disjoint and $x$ and $z$ are disjoint, then $x$ and
$y \vee z$ are disjoint.
\end{lemma}

\begin{lemma}
\label{lem:isVLDisjoint-add-right}
\lean{IsVLDisjoint.add_right}
\leanok
\uses{def:isVLDisjoint, lem:inf-le-inf-add-inf-of-nonneg}
Let $X$ be a lattice-ordered abelian group and let $x, y, z \in X$. If
$x$ and $y$ are disjoint and $x$ and $z$ are disjoint, then $x$ and
$y + z$ are disjoint.
\end{lemma}

\begin{lemma}
\label{lem:isVLDisjoint-sup-left}
\lean{IsVLDisjoint.sup_left}
\leanok
\uses{lem:isVLDisjoint-sup-right, lem:isVLDisjoint-comm}
Let $X$ be a lattice-ordered abelian group and let $x, y, z \in X$. If
$x$ and $z$ are disjoint and $y$ and $z$ are disjoint, then $x \vee y$
and $z$ are disjoint.
\end{lemma}

\begin{lemma}
\label{lem:isVLDisjoint-add-left}
\lean{IsVLDisjoint.add_left}
\leanok
\uses{lem:isVLDisjoint-add-right, lem:isVLDisjoint-comm}
Let $X$ be a lattice-ordered abelian group and let $x, y, z \in X$. If
$x$ and $z$ are disjoint and $y$ and $z$ are disjoint, then $x + y$ and
$z$ are disjoint.
\end{lemma}

\begin{lemma}
\label{lem:isVLDisjoint-sub-inf}
\lean{isVLDisjoint_sub_inf}
\leanok
\uses{def:isVLDisjoint, lem:sub-inf-eq-posPart, lem:isVLDisjoint-posPart-negPart}
Let $X$ be a lattice-ordered abelian group and let $x, y \in X$. Then
$x - x \wedge y$ and $y - x \wedge y$ are disjoint.
\end{lemma}

\begin{theorem}
\label{thm:isVLDisjoint-posPart-sub-smul}
\lean{isVLDisjoint_posPart_sub_smul}
\leanok
\uses{def:isVLDisjoint, lem:isVLDisjoint-posPart-negPart, lem:isVLDisjoint-smul-right, lem:nonneg-smul-sup}
Let $X$ be a vector lattice, let $x, y \in X$ and $\lambda \in \mathbb{R}$
with $0 < \lambda$. Then $(x - \lambda y)^+$ and $(y - \lambda^{-1} x)^+$
are disjoint.
\end{theorem}

\begin{theorem}
\label{thm:sum-eq-sup-of-pairwise-isVLDisjoint-of-nonneg}
\lean{sum_eq_sup_of_pairwise_isVLDisjoint_of_nonneg}
\leanok
\uses{def:isVLDisjoint, lem:add-eq-sup-of-isVLDisjoint-of-nonneg, lem:isVLDisjoint-add-right}
Let $X$ be a lattice-ordered abelian group, let $n \in \mathbb{N}$ with
$n > 0$, and let $f : \{0, \dots, n-1\} \to X$ be a family of
non-negative elements which is pairwise disjoint. Then
$\sum_i f_i = \bigvee_i f_i$.
\end{theorem}

\begin{theorem}
\label{thm:abs-sum-of-pairwise-isVLDisjoint}
\lean{abs_sum_of_pairwise_isVLDisjoint}
\leanok
\uses{def:isVLDisjoint, thm:abs-add-of-isVLDisjoint, thm:abs-smul, lem:isVLDisjoint-smul-left, lem:isVLDisjoint-smul-right, lem:isVLDisjoint-add-right}
Let $X$ be a vector lattice, let $n \in \mathbb{N}$, and let
$x : \{0, \dots, n-1\} \to X$ be a pairwise-disjoint family. For every
$\alpha : \{0, \dots, n-1\} \to \mathbb{R}$,
$\lvert \sum_i \alpha_i x_i \rvert = \sum_i \lvert \alpha_i \rvert
\lvert x_i \rvert$.
\end{theorem}

\begin{theorem}
\label{thm:sup-smul-eq-sum-of-pairwise-isVLDisjoint-of-nonneg}
\lean{sup_smul_eq_sum_of_pairwise_isVLDisjoint_of_nonneg}
\leanok
\uses{thm:sum-eq-sup-of-pairwise-isVLDisjoint-of-nonneg, lem:isVLDisjoint-smul-left, lem:isVLDisjoint-smul-right}
Let $X$ be a vector lattice, let $n \in \mathbb{N}$ with $n > 0$, and let
$x : \{0, \dots, n-1\} \to X$ be a pairwise-disjoint family of
non-negative elements. For every non-negative family
$\alpha : \{0, \dots, n-1\} \to \mathbb{R}$,
$\sum_i \alpha_i x_i = \bigvee_i \alpha_i x_i$.
\end{theorem}

\begin{theorem}
\label{thm:linearIndependent-of-pairwise-isVLDisjoint}
\lean{linearIndependent_of_pairwise_isVLDisjoint}
\leanok
\uses{thm:abs-sum-of-pairwise-isVLDisjoint, lem:abs-eq-zero-iff-zero}
Let $X$ be a vector lattice, let $n \in \mathbb{N}$, and let
$x : \{0, \dots, n-1\} \to X$ be a pairwise-disjoint family of non-zero
elements. Then $x$ is linearly independent over $\mathbb{R}$.
\end{theorem}

\begin{theorem}
\label{thm:eq-zero-of-pairwise-isVLDisjoint-tendsto}
\lean{eq_zero_of_pairwise_isVLDisjoint_tendsto}
\leanok
\uses{def:isVLDisjoint, def:normed-vector-lattice, thm:abs-add-of-isVLDisjoint}
Let $Y$ be a normed vector lattice, let $u : \mathbb{N} \to Y$ be a
pairwise-disjoint sequence, and let $x \in Y$. If $u_n \to x$ in norm,
then $x = 0$.
\end{theorem}

\section{Order completeness}

A lattice is \emph{Dedekind complete} (or \emph{order complete}) when
every non-empty subset that is bounded above admits a supremum, and its
\emph{sigma} analogue is obtained by restricting the requirement to
countable subsets. On a vector lattice either property can be checked
on the restricted class of non-negative elements, and indeed only on
increasing bounded above nets (respectively sequences) of non-negative
elements. The section closes by showing that every Dedekind complete
lattice is sigma Dedekind complete, and that sigma Dedekind
completeness already forces the Archimedean property in any
lattice-ordered abelian group.

On a vector lattice, Dedekind completeness can be checked either on
increasing bounded above nets of non-negative elements, or on
non-empty bounded above sets of non-negative elements.

\begin{theorem}
\label{thm:conditionallyCompleteLatticeOfPosNet}
\lean{conditionallyCompleteLatticeOfPosNet}
\leanok
\uses{def:vector-lattice}
Let $X$ be a vector lattice. If every bounded above increasing net of
non-negative elements of $X$ admits a least upper bound, then $X$ is
Dedekind complete.
\end{theorem}

\begin{theorem}
\label{thm:conditionallyCompleteLatticeOfPosSet}
\lean{conditionallyCompleteLatticeOfPosSet}
\leanok
\uses{def:vector-lattice}
Let $X$ be a vector lattice. If every non-empty bounded above set of
non-negative elements of $X$ admits a least upper bound, then $X$ is
Dedekind complete.
\end{theorem}

\subsection{Sigma order completeness}

A lattice is \emph{sigma Dedekind complete} (or \emph{sigma order
complete}) when every non-empty countable subset that is bounded above
admits a supremum, and every non-empty countable subset that is
bounded below admits an infimum.

\begin{definition}
\label{def:sigmaConditionallyCompleteLattice}
\lean{SigmaConditionallyCompleteLattice}
\leanok
A lattice $X$ is \emph{sigma Dedekind complete} when every non-empty
countable subset of $X$ which is bounded above admits a supremum, and
every non-empty countable subset of $X$ which is bounded below admits
an infimum.
\end{definition}

\begin{lemma}
\label{lem:conditionallyCompleteLattice-toSigmaConditionallyCompleteLattice}
\lean{ConditionallyCompleteLattice.toSigmaConditionallyCompleteLattice}
\leanok
\uses{def:sigmaConditionallyCompleteLattice}
Every Dedekind complete lattice is sigma Dedekind complete.
\end{lemma}

As with Dedekind completeness, sigma Dedekind completeness on a
vector lattice can be checked on the restricted class of non-negative
elements, or even only on increasing bounded above sequences of
non-negative elements.

\begin{theorem}
\label{thm:sigmaConditionallyCompleteLatticeOfPosSeq}
\lean{sigmaConditionallyCompleteLatticeOfPosSeq}
\leanok
\uses{def:vector-lattice, def:sigmaConditionallyCompleteLattice}
Let $X$ be a vector lattice. If every bounded above increasing
sequence of non-negative elements of $X$ admits a least upper bound,
then $X$ is sigma Dedekind complete.
\end{theorem}

\begin{theorem}
\label{thm:sigmaConditionallyCompleteLatticeOfPosCountableSet}
\lean{sigmaConditionallyCompleteLatticeOfPosCountableSet}
\leanok
\uses{def:vector-lattice, def:sigmaConditionallyCompleteLattice}
Let $X$ be a vector lattice. If every non-empty bounded above
countable set of non-negative elements of $X$ admits a least upper
bound, then $X$ is sigma Dedekind complete.
\end{theorem}

Sigma conditional completeness alone forces the Archimedean
property in the lattice-ordered group setting.

\begin{theorem}
\label{thm:isVLArchimedean-of-sigmaConditionallyCompleteLattice}
\lean{IsVLArchimedean_of_sigmaConditionallyCompleteLattice}
\leanok
\uses{def:is-vl-archimedean, def:sigmaConditionallyCompleteLattice}
Every sigma Dedekind complete lattice-ordered abelian group is
Archimedean.
\end{theorem}

